\section{Proof of the Affine Delayed Markov Bound}
Fix a stationary action $a=(q,b,\eta)$ and condition on the history immediately before the retained joint observation.
All expectations below include both the Markov observation and the random delay profile.
Write $\bar H_k=q^{-1}\sum_iH_{i,k}$, $\bar\xi_k=q^{-1}\sum_i\xi_{i,k}$, and
\begin{equation}
v_k=e_k-\eta(\bar H_ke_k-\bar\xi_k),
\quad
r_k=\frac1q\sum_iH_{i,k}(e_{k-\tau_i}-e_k).
\label{eq:appendix-vr}
\end{equation}
Then \eqref{eq:affine-td} is $e_{k+1}=v_k-\eta r_k$.

Total-variation duality and stationary centering give
\begin{equation}
\left\|\E[\bar\xi_k\mid\F_k]\right\|\le2G\delta.
\label{eq:conditional-bias}
\end{equation}
Moreover, stationarity and the bounded integrand imply
\begin{align}
\E[\|\bar H_ke-\bar\xi_k\|^2\mid\F_k]
&\le(2K_q+4L^2\delta)\|e\|^2\nonumber\\
&\quad+2\Omega_q+4G^2\delta.
\label{eq:second-moment-transfer}
\end{align}
Combining \eqref{eq:conditional-bias} with the strong monotonicity of $A$ and applying Young's inequality with parameter $\mu_\delta$ yields
\begin{equation}
\E\|v_k\|^2\le a_\delta(\eta)R_k+\beta_\delta(\eta).
\label{eq:v-bound}
\end{equation}

The affine telescoping identity retains the innovation:
\begin{equation}
e_{k-\tau_i}-e_k=\eta\sum_{s=k-\tau_i}^{k-1}\left[\frac1q\sum_jH_{j,s}e_{s-\tau_j}-\bar\xi_s\right].
\end{equation}
Jensen's inequality and the almost-sure bounds in Assumption~\ref{ass:catalogue} give
\begin{equation}
\E\|r_k\|^2\le2\eta^2L^4\tau_{\rm rms}^2R_k+2\eta^2L^2G^2\tau_{\rm rms}^2.
\label{eq:r-bound}
\end{equation}
For any $\lambda>0$,
\begin{equation}
\|v_k-\eta r_k\|^2\le(1+\lambda)\|v_k\|^2+(1+\lambda^{-1})\eta^2\|r_k\|^2.
\end{equation}
Substituting \eqref{eq:v-bound} and \eqref{eq:r-bound} gives $\E\|e_{k+1}\|^2\le c_\lambda R_k+d_\lambda$.
The choice $\lambda=\sqrt{h_\tau/a_\delta}$ minimizes the contraction coefficient and produces \eqref{eq:affine-cd}.
For zero delay, $r_k=0$ and the continuous limiting definitions follow directly.

Let $R_\star=d_{\rm aff}/(1-c_{\rm aff})$.
If $R_k\ge R_\star$, every new value entering the sliding envelope is at most $c_{\rm aff}R_k+d_{\rm aff}\le R_k$.
After $2D+1$ updates, every value in the old envelope has been replaced, so $R_{k+2D+1}\le c_{\rm aff}R_k+d_{\rm aff}$.
If $R_k<R_\star$, the same inequality keeps the enlarged envelope below $R_\star$.
Iterating the two cases proves Theorem~\ref{thm:affine}.

\subsection{Return transfer and arbitrary terminal indices}
Let $n=\lfloor N_a/(2D+1)\rfloor$.
The envelope argument used above shows that all iterates after index $n(2D+1)$ and before the next completed replacement block remain below the larger of the current envelope and $R_\star$.
The right-hand side of \eqref{eq:affine-bound} already has this form because its transient part is truncated by $(R_0-R_\star)_+$ in \eqref{eq:certificate-score}.

Furthermore,
\begin{equation}
J_\nu(w)-J_\nu(w^\star)=\psi_\nu^\top e,
\quad
\psi_\nu=\E_{s\sim\nu}\phi(s).
\end{equation}
Cauchy--Schwarz gives $|\psi_\nu^\top e|^2\le\|\psi_\nu\|^2\|e\|^2$.
Taking expectations and applying the terminal envelope proves Corollary~\ref{cor:return}.

\subsection{Proof of Theorem~\ref{thm:tracking}}
During epoch $r$, hold the target $w_r^\star$ fixed and apply one replacement-block step from the proof of Theorem~\ref{thm:affine}.
If $\widetilde Z_{r+1}$ denotes the terminal history envelope measured relative to $w_r^\star$, then
\begin{equation}
\widetilde Z_{r+1}\le c_{\rm aff}(a_r)Z_r+d_{\rm aff}(a_r)\le cZ_r+d.
\label{eq:tracking-old-target}
\end{equation}
For every iterate $w_s$ in that terminal history and every $\chi>0$, Young's inequality gives
\begin{align}
\|w_s-w_{r+1}^\star\|^2
&\le(1+\chi)\|w_s-w_r^\star\|^2\nonumber\\
&\quad +(1+\chi^{-1})\|w_{r+1}^\star-w_r^\star\|^2.
\end{align}
Taking expectations and the maximum over the history yields
\begin{equation}
Z_{r+1}\le (1+\chi)cZ_r+(1+\chi)d+(1+\chi^{-1})v^2.
\end{equation}
Iteration proves \eqref{eq:tracking-bound} whenever $\bar c=(1+\chi)c<1$.

For a one-step TD advantage,
\begin{equation}
\widehat A_w-A^\star
=\{\gamma\phi(s')-\phi(s)\}^\top(w-w_n^\star).
\end{equation}
The feature bound implies
$\|\gamma\phi(s')-\phi(s)\|\le(1+\gamma)\Phi$.
Cauchy--Schwarz followed by the definition of $Z_n$ proves \eqref{eq:advantage-transfer}.

\subsection{Existence and complexity of the mixed controller}
For fixed $(q,b)$, $a_\delta$, $\beta_\delta$, $h_\tau$, and $g_\tau$ are continuous functions of $\eta$.
The optimized Young coefficient and hence $c_{\rm aff}$ and $d_{\rm aff}$ are continuous wherever their denominators are positive, with the zero-delay values defined by continuity.
The score in \eqref{eq:certificate-score} is therefore continuous on each certified component.
Restricting to its closed feasible portion gives existence by the extreme-value theorem.

The controller evaluates at most $N_qN_b$ scalar problems.
If a deterministic scalar routine uses $L_\eta$ score evaluations per pair, the control arithmetic is $O(N_qN_bL_\eta)$.
For rank-one $H_{i,k}$, matrix--vector products can be accumulated from feature inner products without forming a dense $d\times d$ matrix.
The calibration update is consequently $O(qd)$ in arithmetic and $O(d)$ in working memory, which proves Proposition~\ref{prop:computation}.

\section{Correlation-Limited Minimax Risk}
For one round of the Gaussian location subclass, the agent average has variance $\sigma_c^2+\sigma_e^2/q$.
The $T$ round averages are independent and Gaussian with known variance, so their sample mean is minimax for squared loss over the unrestricted location family and has risk $T^{-1}(\sigma_c^2+\sigma_e^2/q)$.
Dividing the single-agent risk by this value gives \eqref{eq:speedup}.
Since $1+(q-1)\rho\ge q\rho$ for $\rho>0$, the speedup is at most $\rho^{-1}$.

\section{Proof of the Stopped Adaptive Information Identity}
Condition on the pre-action history and realized action $A_k$.
An orthogonal rotation maps $X_k$ to the normalized mean $Y_k=q_k^{-1/2}\bm{1}^\top X_k$ and $q_k-1$ contrasts.
The contrasts are independent standard normal variables under both hypotheses and therefore contribute zero to the likelihood ratio.

The Kalman predictive law of $Y_k$ under hypothesis $j$ is $\mathcal{N}(\sqrt{q_k}m^-_{j,k},V_{j,k})$.
Factoring the stopped path law into action kernels and conditional observation kernels makes the common predictable action factors cancel pointwise.
The log likelihood ratio is therefore the sum of scalar innovation log likelihood ratios.
Conditional expectation under hypothesis zero gives the Gaussian divergence in \eqref{eq:conditional-kl}.
Tonelli's theorem applies because each conditional divergence is nonnegative and the budget bounds $\tau$.
This proves \eqref{eq:adaptive-kl}; exchanging the hypotheses proves the reverse identity.

For a decision $\widehat J$, data processing gives
\begin{equation}
\KL(P_j^\pi|\F_\tau\,\|\,P_{1-j}^\pi|\F_\tau)\ge\operatorname{kl}(\Pp_j\{\widehat J=j\},\Pp_{1-j}\{\widehat J=j\}).
\end{equation}
If both directional errors are at most $\delta$, monotonicity of binary relative entropy yields \eqref{eq:info-constraint}.

\section{Occupation-Measure Opportunity Lower Bound}
Let $\Delta_j(z,a)\ge0$ be the one-step opportunity loss relative to the instance-$j$ certified oracle.
For a stopped predictable policy, define
\begin{equation}
\nu_j(G,a)=\E_j^\pi\sum_{k\le\tau}\mathbf{1}\{Z_k\in G,A_k=a\}.
\end{equation}
Every $\delta$-correct policy induces an occupation measure satisfying the controlled Kalman flow constraints,
\begin{align}
\int\mathcal{I}_{j,1-j}(z,a)\nu_j(dz,da)&\ge\operatorname{kl}(1-\delta,\delta),\nonumber\\
\int(h+q(a))\nu_j(dz,da)&\le B_{\rm msg},\nonumber\\
\int b(a)\nu_j(dz,da)+D\Pp_j(\tau>0)&\le B_{\rm env}.
\end{align}
Hence its expected identification opportunity cost is at least the infimum of $\int\Delta_j(z,a)\nu_j(dz,da)$ over this feasible occupation set.
This formulation retains the belief-state dependence created by irregular Markov gaps and changing participation.

For completeness, let $T_a(z,dz')$ be the belief-state transition kernel created by action $a$ and the scalar Kalman innovation.
The stopped occupation measure and terminal belief law $\zeta_j$ obey, for every bounded measurable test function $f$,
\begin{align}
\int f(z)\zeta_j(dz)-f(z_1)
&=\int\left[\int f(z')T_a(z,dz')\right.\nonumber\\
&\left.\hspace{18mm}-f(z)\right]\nu_j(dz,da).
\label{eq:belief-flow}
\end{align}
This is the controlled flow constraint included in ${\cal M}_j$.
Predictability ensures that $T_a$ is chosen before the current innovation is observed.

The expected accumulated opportunity loss is exactly
\begin{equation}
\E_j^\pi\sum_{k\le\tau}\Delta_j(Z_k,A_k)
=\int\Delta_j(z,a)\nu_j(dz,da).
\end{equation}
Equations \eqref{eq:adaptive-kl}, \eqref{eq:budgets}, and \eqref{eq:belief-flow} show that the occupation measure of every admissible policy belongs to ${\cal M}_j$.
Taking an infimum over the larger set of all feasible measures gives \eqref{eq:occupation-lower-bound} and proves Proposition~\ref{prop:occupation}.

\section{Proof of the Main Learning-Value Guarantee}
Condition on whether the likelihood test identifies instance $j$ correctly.
On the correct event, the post-probe risk is at most $\overline R_j^\star(B_d^D,D)$.
On the wrong event, its excess over the correct catalogue optimum is at most $W_j(d)$.
Since the wrong-event probability is at most $\delta$, total expectation proves \eqref{eq:main-risk}.

Subtracting the full-budget all-agent risk and using \eqref{eq:value-decomposition} gives
\begin{equation}
\E_1\|e_{\rm LA}\|^2-\overline R_1^{\rm all}(B,D)
\le C_1^{\rm probe}+C_1^D-G_1+\delta W_1<0.
\end{equation}
The same decomposition on instance zero is bounded above by $\epsilon_{\rm safe}(d)\le\bar\epsilon$.
The fallback statement follows directly from the strict admission rule.

\section{Proof of the Threshold Sandwich}
Write $d_\delta=\operatorname{kl}(1-\delta,\delta)$.
Compact separation, $\lambda\le1-\gamma$, and continuity of the scalar Kalman recursion imply that the directional information of every feasible observation is at most a finite class constant $i_+$.
Every action has a positive resource cost, so the information constraint \eqref{eq:info-constraint} implies $B_N\ge c_Nd_\delta$ for a constant $c_N>0$ depending only on the class and the budget ray.

For the fixed probe in the theorem assumptions, compactness gives a positive lower Bhattacharyya information $i_-$ per retained observation.
Thus a block of
\begin{equation}
n_\delta=\left\lceil i_-^{-1}\log(1/\delta)\right\rceil
\end{equation}
has both directional errors at most $\delta$.
Its two resource costs are at most $c_pn_\delta+D_{\max}$ in budget-ray units.
For $0<\delta\le\delta_0<1/2$, continuity of $\log(1/\delta)/d_\delta$ on $(0,\delta_0]$, together with its finite limit at zero, yields
\begin{equation}
n_\delta\le c_0+c_1d_\delta\le c_0+c_2B_N.
\label{eq:probe-versus-necessary}
\end{equation}

For a stationary action $a=(q,b,\eta)$, the exact running-mean risk has the expansion
\begin{equation}
R_j^{\rm mean}(a;s)=\frac{K_j(a)}{s}+O_{\mathcal K}(s^{-2})
\end{equation}
uniformly along the budget ray.
The same expansion is assumed for every certified catalogue risk.
Removing $n_\delta$ probe blocks and at most $D_{\max}$ terminal transitions perturbs each certificate by at most
\begin{equation}
\frac{L_{\mathcal K}(n_\delta+D_{\max}+1)}{s^2}
\label{eq:risk-perturbation}
\end{equation}
for a class constant $L_{\mathcal K}$ once all catalogue actions are feasible.
The high-instance all-agent-to-oracle improvement is at least $g_{\min}/s-O_{\mathcal K}(s^{-2})$.
By the wrong-commit assumption, the term $\delta W_1$ consumes at most half of this leading gain.
Equations \eqref{eq:probe-versus-necessary}--\eqref{eq:risk-perturbation} therefore make the strict high-instance admission inequality hold when
\begin{equation}
s\ge c_3B_N+c_4(B_{\rm oracle}+1).
\end{equation}
All low-instance certificate differences are $O_{\mathcal K}(s^{-1})$.
For a fixed admissible safety target, they fall below that target once $s\ge B_\epsilon$.
Taking the maximum of these sufficient scales and increasing class constants gives
\begin{equation}
B_S\le C_{\mathcal K}B_N+C'_{\mathcal K}(B_{\rm oracle}+B_\epsilon+1).
\end{equation}
The inequality $B_N\le B_S$ follows from the necessity of the two directional information constraints for every admitted $\delta$-correct probe.
This proves \eqref{eq:threshold-sandwich}.

\section{Secondary SDDE Interpretation}
The controlled discrete recursion has the local stochastic delay differential equation representation
\begin{equation}
dE(t)=-\frac1q\sum_{i=1}^{q}A_iE(t-\tau_i(t))dt+\sqrt{\eta}\,\Sigma_q^{1/2}(E_t)dW(t).
\end{equation}
The diffusion must retain the state-dependent quadratic variation induced by random temporal-difference Jacobians.
Near the fixed point, its multiplicative component obeys a quadratic bound proportional to $K_q\|E_t\|^2$.
This representation explains why correlation changes both the noise floor and the mean-square stability margin.
The paper's finite-time guarantees are provided by the exact discrete-time Lyapunov analysis above.
